\section{Derived categories}
\label{section-derived-categories}

\noindent
In this section we construct the derived category of an abelian category
$\mathcal{A}$ by inverting the quasi-isomorphisms in $K(\mathcal{A})$.
Before we do this recall that the functors
$H^i : \text{Comp}(\mathcal{A}) \to \mathcal{A}$
factor through $K(\mathcal{A})$, see
Homology, Lemma \ref{homology-lemma-map-cohomology-homotopy-cochain}.
Moreover, in
Homology, Definition \ref{homology-definition-cohomology-shift}
we have defined identifications $H^i(K^\bullet[n]) = H^{i + n}(K^\bullet)$.
At this point it makes sense to redefine
$$
H^i(K^\bullet) = H^0(K^\bullet[i])
$$
in order to avoid confusion and possible sign errors.

\begin{lemma}
\label{lemma-cohomology-homological}
Let $\mathcal{A}$ be an abelian category. The functor
$$
H^0 : K(\mathcal{A}) \longrightarrow \mathcal{A}
$$
is homological.
\end{lemma}

\begin{proof}
Because $H^0$ is a functor, and by our definition of distinguished triangles
it suffices to prove that given a termwise split short exact sequence
of complexes $0 \to A^\bullet \to B^\bullet \to C^\bullet \to 0$
the sequence $H^0(A^\bullet) \to H^0(B^\bullet) \to H^0(C^\bullet)$
is exact. This follows from
Homology, Lemma \ref{homology-lemma-long-exact-sequence-cochain}.
\end{proof}

\noindent
In particular, this lemma implies that a distinguished triangle
$(X, Y, Z, f, g, h)$ in $K(\mathcal{A})$ gives rise to a long exact
cohomology sequence
\begin{equation}
\label{equation-long-exact-cohomology-sequence-D}
\xymatrix{
\ldots \ar[r] &
H^i(X) \ar[r]^{H^i(f)} &
H^i(Y) \ar[r]^{H^i(g)} &
H^i(Z) \ar[r]^{H^i(h)} &
H^{i + 1}(X) \ar[r] & \ldots
}
\end{equation}
see (\ref{equation-long-exact-cohomology-sequence}). Moreover, there is
a compatibility with the long exact sequence of cohomology associated to
a short exact sequence of complexes (insert future reference here). For
example, if $(A^\bullet, B^\bullet, C^\bullet, \alpha, \beta, \delta)$
is the distinguished triangle associated to a termwise split exact
sequence of complexes (see
Definition \ref{definition-split-ses}),
then the cohomology sequence above agrees with the one defined using the
snake lemma, see
Homology, Lemma \ref{homology-lemma-long-exact-sequence-cochain}
and for agreement of sequences, see
Homology, Lemma \ref{homology-lemma-ses-termwise-split-long-cochain}.

\medskip\noindent
Recall that a complex $K^\bullet$ is {\it acyclic} if $H^i(K^\bullet) = 0$
for all $i \in \mathbf{Z}$. Moreover, recall that a morphism of complexes
$f : K^\bullet \to L^\bullet$ is a {\it quasi-isomorphism} if and only if
$H^i(f)$ is an isomorphism for all $i$. See
Homology, Definition \ref{homology-definition-quasi-isomorphism-cochain}.

\begin{lemma}
\label{lemma-acyclic}
Let $\mathcal{A}$ be an abelian category. The full subcategory
$\text{Ac}(\mathcal{A})$ of $K(\mathcal{A})$ consisting of acyclic complexes
is a strictly full saturated triangulated subcategory of $K(\mathcal{A})$.
The corresponding saturated multiplicative system (see
Lemma \ref{lemma-operations})
of $K(\mathcal{A})$ is the set $\text{Qis}(\mathcal{A})$
of quasi-isomorphisms. In particular, the kernel of the localization
functor $Q : K(\mathcal{A}) \to \text{Qis}(\mathcal{A})^{-1}K(\mathcal{A})$
is $\text{Ac}(\mathcal{A})$ and the functor $H^0$ factors through $Q$.
\end{lemma}

\begin{proof}
We know that $H^0$ is a homological functor by
Lemma \ref{lemma-cohomology-homological}.
Thus this lemma is a special case of
Lemma \ref{lemma-acyclic-general}.
\end{proof}

\begin{definition}
\label{definition-unbounded-derived-category}
Let $\mathcal{A}$ be an abelian category.
Let $\text{Ac}(\mathcal{A})$ and $\text{Qis}(\mathcal{A})$
be as in
Lemma \ref{lemma-acyclic}.
The {\it derived category of $\mathcal{A}$} is the triangulated
category
$$
D(\mathcal{A}) =
K(\mathcal{A})/\text{Ac}(\mathcal{A}) =
\text{Qis}(\mathcal{A})^{-1} K(\mathcal{A}).
$$
We denote $H^0 : D(\mathcal{A}) \to \mathcal{A}$ the unique functor
whose composition with the quotient functor gives back the functor
$H^0$ defined above. Using
Lemma \ref{lemma-homological-functor-bounded}
we introduce the strictly full saturated triangulated subcategories
$D^{+}(\mathcal{A}), D^{-}(\mathcal{A}), D^b(\mathcal{A})$
whose sets of objects are
$$
\begin{matrix}
\Ob(D^{+}(\mathcal{A})) =
\{X \in \Ob(D(\mathcal{A})) \mid
H^n(X) = 0\text{ for all }n \ll 0\} \\
\Ob(D^{-}(\mathcal{A})) =
\{X \in \Ob(D(\mathcal{A})) \mid
H^n(X) = 0\text{ for all }n \gg 0\} \\
\Ob(D^b(\mathcal{A})) =
\{X \in \Ob(D(\mathcal{A})) \mid
H^n(X) = 0\text{ for all }|n| \gg 0\}
\end{matrix}
$$
The category $D^b(\mathcal{A})$ is called the {\it bounded derived
category} of $\mathcal{A}$.
\end{definition}

\noindent
If $K^\bullet$ and $L^\bullet$ are complexes of $\mathcal{A}$
then we sometimes say ``$K^\bullet$ is {\it quasi-isomorphic} to
$L^\bullet$'' to indicate that $K^\bullet$ and $L^\bullet$ are
isomorphic objects of $D(\mathcal{A})$.

\begin{remark}
\label{remark-existence-derived}
In this chapter, we consistently work with ``small'' abelian categories
(as is the convention in the Stacks project). For a ``big'' abelian
category $\mathcal{A}$, it isn't clear that the derived category
$D(\mathcal{A})$ exists, because it isn't clear that morphisms in the
derived category are sets. In fact, in general they aren't, see
Examples, Lemma \ref{examples-lemma-big-abelian-category}.
However, if $\mathcal{A}$ is a Grothendieck abelian category, and given
$K^\bullet, L^\bullet$ in $K(\mathcal{A})$, then by
Injectives, Theorem \ref{injectives-theorem-K-injective-embedding-grothendieck}
there exists a quasi-isomorphism $L^\bullet \to I^\bullet$ to a
K-injective complex $I^\bullet$ and Lemma \ref{lemma-K-injective} shows that
$$
\Hom_{D(\mathcal{A})}(K^\bullet, L^\bullet) =
\Hom_{K(\mathcal{A})}(K^\bullet, I^\bullet)
$$
which is a set. Some examples of Grothendieck abelian categories
are the category of modules over a ring, or more generally
the category of sheaves of modules on a ringed site.
\end{remark}

\noindent
Each of the variants $D^{+}(\mathcal{A}), D^{-}(\mathcal{A}), D^b(\mathcal{A})$
can be constructed as a localization of the corresponding homotopy category.
This relies on the following simple lemma.

\begin{lemma}
\label{lemma-complex-cohomology-bounded}
Let $\mathcal{A}$ be an abelian category.
Let $K^\bullet$ be a complex.
\begin{enumerate}
\item If $H^n(K^\bullet) = 0$ for all $n \ll 0$, then there exists
a quasi-isomorphism $K^\bullet \to L^\bullet$ with $L^\bullet$
bounded below.
\item If $H^n(K^\bullet) = 0$ for all $n \gg 0$, then there exists
a quasi-isomorphism $M^\bullet \to K^\bullet$ with $M^\bullet$
bounded above.
\item If $H^n(K^\bullet) = 0$ for all $|n| \gg 0$, then there exists
a commutative diagram of morphisms of complexes
$$
\xymatrix{
K^\bullet \ar[r] & L^\bullet \\
M^\bullet \ar[u] \ar[r] & N^\bullet \ar[u]
}
$$
where all the arrows are quasi-isomorphisms, $L^\bullet$
bounded below, $M^\bullet$ bounded above, and $N^\bullet$ a bounded
complex.
\end{enumerate}
\end{lemma}

\begin{proof}
Pick $a \ll 0 \ll b$ and set $M^\bullet = \tau_{\leq b}K^\bullet$,
$L^\bullet = \tau_{\geq a}K^\bullet$, and
$N^\bullet = \tau_{\leq b}L^\bullet = \tau_{\geq a}M^\bullet$.
See
Homology, Section \ref{homology-section-truncations}
for the truncation functors.
\end{proof}

\noindent
To state the following lemma denote
$\text{Ac}^{+}(\mathcal{A})$, $\text{Ac}^{-}(\mathcal{A})$,
resp.\ $\text{Ac}^b(\mathcal{A})$ the intersection of
$K^{+}(\mathcal{A})$, $K^{-}(\mathcal{A})$, resp.\ $K^b(\mathcal{A})$
with $\text{Ac}(\mathcal{A})$.
Denote $\text{Qis}^{+}(\mathcal{A})$, $\text{Qis}^{-}(\mathcal{A})$,
resp.\ $\text{Qis}^b(\mathcal{A})$ the intersection of
$K^{+}(\mathcal{A})$, $K^{-}(\mathcal{A})$, resp.\ $K^b(\mathcal{A})$
with $\text{Qis}(\mathcal{A})$.

\begin{lemma}
\label{lemma-bounded-derived}
Let $\mathcal{A}$ be an abelian category. The subcategories
$\text{Ac}^{+}(\mathcal{A})$, $\text{Ac}^{-}(\mathcal{A})$,
resp.\ $\text{Ac}^b(\mathcal{A})$
are strictly full saturated triangulated subcategories
of $K^{+}(\mathcal{A})$, $K^{-}(\mathcal{A})$, resp.\ $K^b(\mathcal{A})$.
The corresponding saturated multiplicative systems (see
Lemma \ref{lemma-operations})
are the sets $\text{Qis}^{+}(\mathcal{A})$, $\text{Qis}^{-}(\mathcal{A})$,
resp.\ $\text{Qis}^b(\mathcal{A})$.
\begin{enumerate}
\item The kernel of the functor $K^{+}(\mathcal{A}) \to D^{+}(\mathcal{A})$
is $\text{Ac}^{+}(\mathcal{A})$ and this induces an equivalence
of triangulated categories
$$
K^{+}(\mathcal{A})/\text{Ac}^{+}(\mathcal{A}) =
\text{Qis}^{+}(\mathcal{A})^{-1}K^{+}(\mathcal{A})
\longrightarrow
D^{+}(\mathcal{A})
$$
\item The kernel of the functor $K^{-}(\mathcal{A}) \to D^{-}(\mathcal{A})$
is $\text{Ac}^{-}(\mathcal{A})$ and this induces an equivalence
of triangulated categories
$$
K^{-}(\mathcal{A})/\text{Ac}^{-}(\mathcal{A}) =
\text{Qis}^{-}(\mathcal{A})^{-1}K^{-}(\mathcal{A})
\longrightarrow
D^{-}(\mathcal{A})
$$
\item The kernel of the functor $K^b(\mathcal{A}) \to D^b(\mathcal{A})$
is $\text{Ac}^b(\mathcal{A})$ and this induces an equivalence
of triangulated categories
$$
K^b(\mathcal{A})/\text{Ac}^b(\mathcal{A}) =
\text{Qis}^b(\mathcal{A})^{-1}K^b(\mathcal{A})
\longrightarrow
D^b(\mathcal{A})
$$
\end{enumerate}
\end{lemma}

\begin{proof}
The initial statements follow from
Lemma \ref{lemma-acyclic-general}
by considering the restriction of the homological functor $H^0$.
The statement on kernels in (1), (2), (3) is a consequence of the
definitions in each case.
Each of the functors is essentially surjective by
Lemma \ref{lemma-complex-cohomology-bounded}.
To finish the proof we have to show the functors are fully faithful.
We first do this for the bounded below version.

\medskip\noindent
Suppose that $K^\bullet, L^\bullet$ are bounded below complexes.
A morphism between these in $D(\mathcal{A})$ is of the form
$s^{-1}f$ for a pair
$f : K^\bullet \to (L')^\bullet$, $s : L^\bullet \to (L')^\bullet$
where $s$ is a quasi-isomorphism. This implies that $(L')^\bullet$
has cohomology bounded below. Hence by
Lemma \ref{lemma-complex-cohomology-bounded}
we can choose a quasi-isomorphism
$s' : (L')^\bullet \to (L'')^\bullet$
with $(L'')^\bullet$ bounded below. Then the pair $(s' \circ f, s' \circ s)$
defines a morphism in $\text{Qis}^{+}(\mathcal{A})^{-1}K^{+}(\mathcal{A})$.
Hence the functor is ``full''. Finally, suppose that the pair
$f : K^\bullet \to (L')^\bullet$, $s : L^\bullet \to (L')^\bullet$
defines a morphism in $\text{Qis}^{+}(\mathcal{A})^{-1}K^{+}(\mathcal{A})$
which is zero in $D(\mathcal{A})$. This means that there exists a
quasi-isomorphism $s' : (L')^\bullet \to (L'')^\bullet$
such that $s' \circ f = 0$. Using
Lemma \ref{lemma-complex-cohomology-bounded}
once more we obtain a quasi-isomorphism
$s'' : (L'')^\bullet \to (L''')^\bullet$
with $(L''')^\bullet$ bounded below.
Thus we see that $s'' \circ s' \circ f = 0$ which implies that
$s^{-1}f$ is zero in $\text{Qis}^{+}(\mathcal{A})^{-1}K^{+}(\mathcal{A})$.
This finishes the proof that the functor in (1) is an equivalence.

\medskip\noindent
The proof of (2) is dual to the proof of (1).
To prove (3) we may use the result of (2). Hence it suffices to
prove that the functor
$\text{Qis}^b(\mathcal{A})^{-1}K^b(\mathcal{A})
\to \text{Qis}^{-}(\mathcal{A})^{-1}K^{-}(\mathcal{A})$
is fully faithful. The argument given in the previous paragraph
applies directly to show this where we consistently work with complexes
which are already bounded above.
\end{proof}







